\section{上下文文件：真正的力量来源}

\begin{importantbox}{关键洞察}
插件只是起点。Cowork 的真正力量来自 \textbf{context files} --- 你创建的 Markdown 文件，告诉 Cowork 你是谁、如何工作。这些文件在每次对话前被读取，让 AI 具备持久记忆。
\end{importantbox}

\subsection{about-me.md --- 个人身份档案}

这是用文本文件表达的职业身份，应包含：

\begin{itemize}
    \item 姓名、角色、公司
    \item 沟通偏好（直接型？详细型？）
    \item 时区和工作时间
    \item 2--3 个写作或工作风格的示例
\end{itemize}

有了这个文件，Cowork 在每次响应前都会读取它，用户不再需要反复解释自己的背景。

\subsection{brand-voice.md --- 品牌语调档案}
\label{sec:brand-voice}

对于任何内容创作者来说，这个文件是必备的：

\begin{itemize}
    \item 你使用的词汇 vs. 你避免的词汇
    \item 语调（casual、professional、bold、supportive 等）
    \item 听起来像你的例句
    \item 你涉及的话题 vs. 你不涉及的话题
\end{itemize}

\begin{knowledgebox}{作者的实践}
作者维护了一个 2000 词的 voice profile 文件。配合 Marketing Plugin 使用后，每篇草稿无需编辑就能呈现其个人风格。这是"context game"的直接体现。
\end{knowledgebox}

\subsection{current-projects.md --- 当前项目档案}

一份反映你当前工作状态的动态文档：

\begin{itemize}
    \item 进行中的项目及其状态
    \item 本周关键 deadline
    \item 阻塞项或待解决问题
    \item 相关文档或资源的链接
\end{itemize}

\begin{warningbox}{需要定期更新}
这个文件必须\textbf{每周更新}。如果文件内容过时，Cowork 基于过期上下文给出的建议反而会产生误导。上下文的价值与其时效性直接相关。
\end{warningbox}

\subsection{本章小结}

三个上下文文件分别解决"我是谁"（about-me.md）、"我怎么说话"（brand-voice.md）、"我在忙什么"（current-projects.md）三个核心问题。它们共同赋予 Cowork 持久化的个人记忆，是从通用 AI 到个人 AI 助手的关键跨越。

